% ============================================================
%  Helge Holst, »Relativitetsprincipet«
%  Fysisk Tidsskrift, 15. Aargang (1916-17), pp. 165-168.
%
%  DIPLOMATIC TRANSCRIPTION (Danish)
%  Status: COMPLETE. Printed pp. 165-168, all 4 pages, transcribed 2026-09-04.
%  No gaps, no duplicates; braces and guillemets balanced; compiles with
%  0 errors, 0 missing characters, 0 overfull boxes. Translation not begun.
%
%  Holst's SECOND contribution to the Kroman relativity debate -- the sixth of
%  its seven items, and short, three pages. Found in the scan 2026-09-04 and
%  recorded in no bibliography consulted. Not to be confused with his
%  »Tidsproblemet« at pp. 1-13 of the same volume, which is
%  ../holst-tidsproblemet. See ../relativitetsprincipet/note.md.
%
%  RIGHTS: clear. Helge Holst died in 1944; Danish copyright (life+70) expired
%  at the end of 2014. Free to transcribe and to publish.
%
% ------------------------------------------------------------
%  THE WITNESS
% ------------------------------------------------------------
%  Google Books `Tr4ZAAAAIAAJ` -- Fysisk tidsskrift, VOLUMES 13-15, the
%  University of California copy, digitised 17 January 2007. Kept as
%  texts/kroman/scan-ft-13-15.pdf, untracked; shared with the siblings.
%
%  CAUTION: this scan's JBIG2 compression drops hairlines -- the fraction rules
%  on ../holst-tidsproblemet's p. 8 are simply absent from the image and had to
%  be confirmed from a parallel on the next page. Check any thin rule at a
%  second resolution before concluding it is not there.
%
% ------------------------------------------------------------
%  PAGE MAP -- VERIFIED 2026-09-04
% ------------------------------------------------------------
%      printed 165-168  =  PDF + 618   (PDF 783-786).
%
%  The offset is NOT the 602 of the front of this volume, nor the 622 of
%  p. 192: unpaginated advertisement leaves are bound in, so vol. 15's offset
%  grows across the volume. Confirmed here on content -- PDF 784 carries the
%  numeral 166, PDF 785 the running head »Relativitetsprincipet.«, PDF 787 the
%  head »Fysisk Tidsskrift« at p. 169.
%
%  CORRECTED from pp. 166-168, which is what the running heads suggest. The
%  article begins in the LOWER HALF of p. 165, under the running head of the
%  article it follows (»Astrospektroskopien og nogle af dens nyere Resultater.
%  II«), with its own title block »Relativitetsprincipet. / Af / Helge Holst.«
%  This is the third article in the debate to do it -- the same thing put
%  ../walsoe-videnskabens-graense at pp. 13-17 rather than 14-17, and Kroman's
%  1917 reply begins partway down p. 192. In this journal, ALWAYS read the page
%  before the apparent first page. It ends a third of the way down p. 168;
%  »Fysik i Skolen« follows.
%
% ------------------------------------------------------------
%  CONVENTIONS
% ------------------------------------------------------------
%  * Diplomatic. 1917 orthography exactly as printed; the pre-1948 aa kept,
%    nouns capitalised, nothing modernised.
%  * QUOTATION MARKS: guillemets »...«, 2 pairs, balanced.
%  * EMPHASIS is LETTERSPACING, set \emph{}, only 8 runs (p. 165 x1, p. 166 x5,
%    p. 167 x2, p. 168 none). PROPER NAMES ARE NEVER SPACED -- Lorentz,
%    Einstein, Michelson all plain -- which is how the compositor sets Holst in
%    »Tidsproblemet« too, and the opposite of how he sets Kroman. No true italic.
%  * FOOTNOTES: exactly one, on p. 166, marked *). The fixed
%    \renewcommand{\thefootnote}{*)} in the preamble is therefore correct here.
%  * MATHEMATICS: one inline $c$ on p. 165. No displays, no fractions, no
%    equation numbers -- so the scan's hairline problem does not arise.
%  * TITLE SPELLING: the title prints »Relativitetsprincipet« with one p, while
%    the body writes »Relativitetsprincippet« with two. Both kept as printed.
%
% ------------------------------------------------------------
%  PRINTER'S DEFECTS -- transcribed AS PRINTED, logged at the site
% ------------------------------------------------------------
%    p. 166  the comma pair round »idet de følger Synkroniseringsreglerne« is
%            not set
%    p. 167  »men at Understregningen« where the sense wants »af«
%  Plate specks, logged as readings rather than defects: »Hånsen« p. 165 and
%  »ęr« p. 167 -- offset ink, not lost or wrong sorts.
%
% ------------------------------------------------------------
%  BUILD
% ------------------------------------------------------------
%  pdflatex path. Sandbox compile-test per TRANSCRIPTION-PLAYBOOK.md §5.
% ============================================================

\documentclass[12pt,a4paper]{article}
\usepackage[utf8]{inputenc}
\usepackage[T1]{fontenc}
\usepackage{libertinus}
\usepackage{libertinust1math}
\usepackage{textalpha}   % Greek in text, typed directly
\usepackage{amsmath}     % this debate has real formulas; see CONVENTIONS
\usepackage{tikz}        % for the redrawn Fizeau diagram, printed p. 8
\usetikzlibrary{arrows.meta}

% FOOTNOTES -- NOT YET ESTABLISHED for this article. The fixed *) below is
% inherited from ../holst-tidsproblemet, where it is correct. It is WRONG
% wherever a page carries two notes (as Walsøe's pp. 13 and 15 do), and it will
% SILENTLY override an \ifcase form placed after it -- that bug cost an hour on
% ../walsoe-videnskabens-graense. Check the pages and replace this line if need be.
\renewcommand{\thefootnote}{*)}
\usepackage[danish]{babel}
\usepackage[protrusion=true,expansion=true]{microtype}
\usepackage{setspace}
\usepackage[margin=1.2in]{geometry}
\usepackage{fancyhdr}
\setlength{\headheight}{15pt}
\usepackage{hyperref}

\hypersetup{
  pdftitle={Relativitetsprincipet (Holst 1917)},
  pdfauthor={Helge Holst},
  hidelinks
}

\pagestyle{fancy}
\fancyhf{}
\fancyhead[LE,RO]{\thepage}
\fancyhead[RE]{\textit{Relativitetsprincipet}}
\fancyhead[LO]{\textit{Helge Holst}}

\setstretch{1.3}

\title{\textbf{\guillemotleft Relativitetsprincipet\guillemotright}}
\author{H.\ Holst}
\date{\textit{Fysisk Tidsskrift} 15 (1916--17), pp.~165--168}

\begin{document}
\maketitle

% --- p. 165 ---
% The article BEGINS ON p. 165, in the lower half of the page, under the
% running head of the preceding article (»Astrospektroskopien og nogle af dens
% nyere Resultater. II«), which ends in the upper half. Title block as printed:
%     Relativitetsprincipet.  /  Af  /  Helge Holst.
% supplied by \maketitle in the preamble and not repeated here.
% NOTE THE SPELLING: the display title and both running heads read
% »Relativitetsprincipet« with ONE p; the first sentence of the body reads
% »Relativitetsprincippet« with TWO. Both are set as printed.
I Anledning af Dr.~H.~M. Hansen's Artikel om Relativitetsprincippet i
% »Hansen« carries a detached accent-shaped ink speck above the a
% (»Hånsen«). Plate dirt, not a sort: read Hansen.
1.~Hefte af denne Aargang vilde jeg gerne bede om Plads for endnu nogle faa
Bemærkninger, i det Haab, at disse kan bidrage til yderligere at klare
Forholdet mellem de forskellige Standpunkter til Einstein's Lære.

For saa vidt man blot anvender den Einstein'ske \emph{Fremgangs\-%
maade} rent fænomenologisk som et Middel til at eliminere tilsyneladende
% the letterspaced run breaks across the printed line: »F r e m g a n g s-«
% / »m a a d e«. Held together with \-%.
Modsigelser (f.~Eks. dem, der blotlagdes gennem det Michelson'ske Forsøg),
kan herimod ikke rejses erkendelsesteoretiske Indvendinger. Man kan
naturligvis vedtage den Regel, at hver enkelt af flere Iagttagere i
forskellige Systemer, der bevæger sig i Forhold til hinanden, lader som om
Lyset i hans System altid har den konstante Hastighed $c$ i alle Retninger
(og at han derfor synkroniserer sine Ure ved Lyssignaler paa den af Einstein
angivne Maade), samt at man gaar over fra det ene af Systemerne til det andet
ved de Einstein'ske Transformationsformler. Man har endvidere Ret til at
opstille den Hypotese, at de Resultater, man finder under Anvendelse af disse
Regler, paa alle Omraader vil stemme med Forsøgsresul\-%
% --- p. 166 ---
taterne. Viser det sig nu, at man ikke blot opnaar saadan Overensstemmelse
ved de Fænomener, for hvilke Fremgangsmaaden er udviklet, men ogsaa for helt
andre Fænomener, ja at man kan forudsige nye Forsøgsresultater, som derefter
virkelig findes, da har Hypotesen hævdet sig som en frugtbar Ledehypotese, og
det er forstaaeligt, at de, der har arbejdet med den, bliver tilbøjelige til
at stole paa dens Rigtighed og anse den som Udtryk for fundamentale
Ejendommeligheder ved Verden. Den Omstændighed, at de forskellige Iagttagere
idet de følger Synkroniseringsreglerne kommer i Uoverensstemmelse ved
% PRINTER'S DEFECT, as printed: the parenthetical »idet de følger
% Synkroniseringsreglerne« is unpunctuated -- no comma after »Iagttagere«
% and none after »Synkroniseringsreglerne«. Verified at 400 dpi; the line
% ends flush at »Iagttagere« with no clipped mark.
Samtidighedsbestemmelser, og at Hypotesen medfører, at der ikke kan findes
noget \emph{Kriterium} paa, hvem der har Ret, gør det i og for sig ikke
erkendelsesteoretisk uantageligt, at Hypotesen om Fremgangsmaadens
Almenanvendelighed er rigtig. Lorentz har jo ved sin
Hypotese\footnote{Baade her og i den tidligere Artikel menes ved Lorentz'
Hypotese eller Fortolkning ikke blot den oprindelige Kontraktionshypotese,
men denne i Forening med dens senere Tilføjelser.} vist, at det (i alt Fald
tilsyneladende) er \emph{mu\-%
ligt} at give en med den naturlige Rum- og Tidsopfattelse fuldstændig
% letterspaced run broken across the printed line: »m u -« / »l i g t«.
forenelig fysisk Forklaring af Sagen.

Men saaledes som Einstein's Lære er fremsat af sin Ophavsmand og almindelig
fremføres, er den ikke blot Beskrivelse af en Fremgangsmaade og Opstilling af
en Hypotese (eller et Dogme) om dennes Almenanvendelighed. Den har foruden en
fænomenologisk Side ogsaa en erkendelsesteoretisk, idet det udtrykkeligt
hævdes eller stiltiende forudsættes, at Reglerne for Samtidighedsbestemmelser
er mere end rent konventionelle, at Lyshastigheden ikke blot \emph{sættes},
men \emph{er} konstant og ens i Forhold til Afsender og Modtager, selv om
disse bevæger sig i Forhold til hinanden, at man ikke kan forbinde nogen
\emph{Mening} med at tale om absolut Samtidighed, og at en fysisk Forklaring
af de omtalte Modsigelser bliver overflødig, naar man revolutionerer sin
Rum- og Tidsopfattelse.

Mod denne Side af Sagen er det, at der har rejst sig Protester. Jeg for mit
Vedkommende har søgt at paavise, at vor Tanke virkelig faar \guillemotleft et
Sted, hvor den kan staa\guillemotright, og hvorfra den kan betragte
Einstein's Lære udefra, naar blot man hævder dens Ret til at forøge
Hastigheder ud over alle Grænser. Hvis vi opgav denne Ret for Tanken, fordi
der (muligvis) er en bestemt fysisk Grænse for opnaaelige Hastigheder, saa
maatte vi f.~Eks. ligesaa vel opgive Tankens Ret til
% --- p. 167 ---
at forlænge rette Linjer ud over alle Grænser, saafremt en eller anden vel
begrundet fysisk Teori havde ført til, at vort Verdenssystem var indesluttet
i en fuldstændig uigennemtrængelig Skal. Heldigvis for Fysikken gør dens
uundværlige Hjælper Matematikken sig ikke nogen fysiske Skrupler af denne
Art; det er en af dens umistelige Rettigheder at lade Tanken gaa ud over det
fysisk opnaaelige. Men naar Tanken ikke lader sig standse af fysiske
Skranker, kan vi, som det er udviklet i min Artikel, forbinde en bestemt
Mening med (og give en Definition af), at to Begivenheder er virkelig
samtidige, selv om Verdensmekanismen skulde være saaledes indrettet, at vi
ikke kan finde noget Kriterium derpaa.
% »Kriterium« is NOT letterspaced here, though it is on p. 166.

Dr.~H. erkender (S.~28, L.~13 f.~n.), at denne Opfattelse af
Samtidighedsbegrebet (hvilken ikke blot falder sammen med den af Lorentz
antydede, men med alle Menneskers stiltiende Forudsætninger --- indtil Aar
1905) kan fastholdes; men at Understregningen af Ordet \emph{kan} samt af
% PRINTER'S DEFECT, as printed: »men at Understregningen ... samt af andre
% Ytringer ... synes det at fremgaa«. The sense wants »men af«. Verified at
% 400 dpi: the sort is a t.
andre Ytringer i Artiklen, synes det at fremgaa, at han mener, det ikke er
nødvendigt for Tanken at fastholde den. Heri er jeg uenig med ham. Den nævnte
Opfattelse kan og maa fastholdes som den, der fremgaar af vore naturlige
Rum- og Tidsbegreber og af vor Tænknings Grundlove. Det er, mener jeg, en
% an ink speck sits under the e of this »er«, giving the look of an ogonek.
% Plate dirt; read »er«.
Illusion, at vi kan komme bort herfra, og det er derfor ogsaa en Illusion, at
Einstein's Lære giver noget som helst Bidrag til \emph{Forklaringen} af de
Michelson'ske Vanskeligheder; den giver blot simple Regler for, hvorledes man
ved Indføring af visse Faktorer kan eliminere de Uoverensstemmelser, som
f.~Eks. opstaar mellem Fysikere, der paa Jorden og paa Mars udfører M.'s
Forsøg.

En anden Sag er det, at naar man anvender Einstein's Fremgangsmaade rent
fænomenologisk, saa gaar hele Spørgsmaalet om Opfattelsen af
% a speck above the g of »Spørgsmaalet« here reads as an acute. Plate dirt.
Samtidighedsbegrebet ud af Diskussionen, og det kan siges at være uden
Interesse for den arbejdende Fysiker, forsaavidt som han foreløbig kan skyde
Sagen til Side.

Utvivlsomt har Dr.~H.'s fyldige og lærerige Redegørelser for Einstein's Lære
været til megen Gavn, og der er Grund til at være ham taknemlig for den
Liberalitet, hvormed han i Tidsskriftet har givet Plads for Indlæg, der gaar
imod hans Synspunkter. Naar han i sin sidste Artikel saa stærkt betoner det
fænomenologiske, aabner han tillige en god Udvej til Udjævning af
Modsætningerne og Udsigt til, at alle Fysikere med Glæde kan følge en
frugtbar Anvendelse af Ein\-%
% --- p. 168 ---
stein's Metode. Dette muliggøres i des højere Grad, jo fuldstændigere det
lykkes at rense Udtryksmaader og Betragtninger for erkendelsesteoretiske
Elementer, der bringer Tanken til at reagere.

Man maa imidlertid haabe, at de Fysikere, der arbejder med Einstein's
Fremgangsmaader, ikke maa være saa fordybede i Arbejdet --- hvor smukke og
store Resultater det end maatte give --- at de helt taber Interessen for
fysiske Forklaringer af Vanskelighederne eller sætter sig saa fast i Troen
paa den Einstein'ske Hypotese, at de paa Forhaand afviser Forklaringsforsøg,
der i sin videre Udformning (i Modsætning til Lorentz') maatte føre til, at
de Einstein'ske Transformationsligninger kun kan give tilnærmelsesvis rigtige
Resultater.
% The article ENDS HERE, about a third of the way down p. 168. No signature
% and no dateline. The rest of p. 168 begins a new item under a new sectional
% head: »Fysik i Skolen.« / »Loven om Aktion og Reaktion i Skolen.« /
% »Af V. A. C. Jensen.«, which runs on to p. 169 and beyond.

\end{document}
